% Generated draft: requires mathematical and source review.
\documentclass[11pt]{article}
\usepackage[margin=25mm]{geometry}
\usepackage{fontspec}
\setmainfont{Latin Modern Roman}
\usepackage{amsmath,amssymb,mathtools,unicode-math}
\setmathfont{Latin Modern Math}
\usepackage{enumitem,needspace,xurl,hyperref}
\setlength{\parindent}{0pt}
\setlength{\parskip}{5pt}
\setlength{\emergencystretch}{3em}
\newcommand{\N}{\mathbb{N}}
\newcommand{\Z}{\mathbb{Z}}
\newcommand{\Q}{\mathbb{Q}}
\newcommand{\R}{\mathbb{R}}
\newcommand{\C}{\mathbb{C}}
\newcommand{\F}{\mathbb{F}}
\newcommand{\A}{\mathbb{A}}
\newcommand{\PP}{\mathbb{P}}
\newcommand{\E}{\mathbb{E}}
\newcommand{\eps}{\varepsilon}
\newcommand{\dd}{\,\mathrm{d}}
\begin{document}
\section*{UnsolvedMath ID: 30001817}
\textbf{Source problem code:} catalogue entry
\subsection{Definitions and mathematical statement}
Let \(p>0\) be a prime, let \(k\) be a perfect field of characteristic \(p\), and let \(W=W(k)\) be its ring of Witt vectors. Let \(X\) be a smooth projective \(W\)-scheme, with special fiber \(X_1=X\times_{\operatorname{Spec}W}\operatorname{Spec}k\) and inclusion \(i:X_1\hookrightarrow X\). For a scheme \(Y\), write \(CH^r(Y)_{\Q}=CH^r(Y)\otimes_{\Z}\Q\) for the group of codimension-\(r\) algebraic cycles modulo rational equivalence with rational coefficients. Here \(r\in\Z_{\geq 0}\) is the conventional codimension range; the supplied statement gives no explicit range for \(r\) and imposes no bound relating \(r\) to \(p\). The restriction \(i^*:CH^r(X)_{\Q}\to CH^r(X_1)_{\Q}\) is the refined Gysin pullback for the special-fiber inclusion.

Write \(H^m_{\mathrm{dR}}(X/W)=\mathbb H^m(X,\Omega^\bullet_{X/W})\), and use a subscript \(\Q\) to denote tensoring over \(\Z\) with \(\Q\). The Hodge filtration is
\[
F^rH^{2r}_{\mathrm{dR}}(X/W)_{\Q}
=\operatorname{im}\!\left(
\mathbb H^{2r}(X,\Omega^{\geq r}_{X/W})\otimes_{\Z}\Q
\longrightarrow H^{2r}_{\mathrm{dR}}(X/W)_{\Q}
\right).
\]
Let
\[
\operatorname{cl}_{\mathrm{cris}}:CH^r(X_1)_{\Q}\longrightarrow H^{2r}_{\mathrm{cris}}(X_1/W)_{\Q}
\]
be the crystalline cycle-class map, and let
\[
c_X:H^{2r}_{\mathrm{cris}}(X_1/W)_{\Q}\xrightarrow{\sim}H^{2r}_{\mathrm{dR}}(X/W)_{\Q}
\]
be the crystalline--de Rham comparison associated with the smooth proper lift \(X\).

\textbf{Question (the supplied p-adic variational Hodge conjecture).} For every such \(p,k,X,r\) and every \(\xi_1\in CH^r(X_1)_{\Q}\), does
\[
c_X\bigl(\operatorname{cl}_{\mathrm{cris}}(\xi_1)\bigr)
\in F^rH^{2r}_{\mathrm{dR}}(X/W)_{\Q}
\]
imply that there exists \(\xi\in CH^r(X)_{\Q}\) satisfying
\[
\operatorname{cl}_{\mathrm{cris}}(i^*\xi)
=\operatorname{cl}_{\mathrm{cris}}(\xi_1)
\quad\text{in }H^{2r}_{\mathrm{cris}}(X_1/W)_{\Q}?
\]
The requested equality is equality of crystalline cycle classes; it does not require \(i^*\xi=\xi_1\) in the Chow group. The supplied original statement phrases this implication declaratively, whereas the normalized statement asks whether it holds. It is retained here as a conjectural question. The supplied partial-solution and verification labels are curation claims, not mathematical conclusions established by this draft.
\subsection{Short English statement}
Given a smooth projective lift over a Witt ring, must every rational cycle on its special fiber whose crystalline class satisfies the indicated Hodge-filtration condition have the same crystalline class as the restriction of a rational cycle on the lift? Equality of the cycles in the Chow group is not requested. The supplied record labels the problem partially solved; this statement-only draft requires human mathematical review.
\subsection{Sources}
\begin{itemize}
\item Local UnsolvedMath registry, supplied record with permanent ID \(30001817\), source problem code \texttt{catalogue entry}, title \emph{The p-Adic Variational Hodge Conjecture}, in the Oberwolfach Reports Open Problems collection. The record gives the proposed year as 2011 and names no individual proposer.
\item \emph{Algebraische Zahlentheorie} (2011), Oberwolfach Reports, pp.\ 1709--1768, DOI \texttt{10.4171/owr/2011/31}: \url{https://doi.org/10.4171/owr/2011/31}. This is the supplied original workshop-report reference; no specific problem page is supplied.
\item Supplied literature source labeled ``Bloch--Esnault--Kerz deformation theorem'': \url{https://arxiv.org/abs/1203.2776}.
\item Supplied literature source labeled ``Langer--motivic Chow groups'': \url{https://doi.org/10.4171/DM/662}.
\item Supplied literature source labeled ``original Oberwolfach report'': \url{https://doi.org/10.4171/owr/2011/31}.
\item Literature assessment dated exactly \texttt{2026-08-21}: supplied dated curation, not independently verified. It labels the problem partially solved and reports that Bloch--Esnault--Kerz prove formal and K-theoretic/vector-bundle versions of the p-adic variational Hodge conjecture, that Langer extends the result to relative motivic Chow groups, and that the full Chow-class lifting statement in the extracted generality is not known. The supplied background describes this as literature triage derived from OWR curation and requests independent verification of time-sensitive status and cited claims.
\item The registry carries the supplied statement-status label \texttt{corrected\_verified} and says that the mathematical content was checked against the cited official report, with notation, context, and question form normalized without changing the source-level claim. That verification is reported from the supplied record and was not independently performed here.
\end{itemize}
\end{document}
